% GENERATED FILE. Edit canonical lessons, then run npm run generate:books.
% canonical-source-hash: fnv1a64:e44d4fdb72504b95
% canonical-lessons: JA-C427-kinchou-suru, JA-C427-dokidoki-suru, JA-C427-sowasowa-suru, JA-C427-iraira-suru, JA-C427-tamerau

\chapter{Nerves Before the Big Day}
\label{ch:ja-nerves-before-the-big-day}
\hlchaptermodality{427}

\begin{chapteropening}
\textbf{By the end of this chapter:} \emph{I can say I am nervous, restless or hesitant.}

Complete the last lesson of chapter 427: I can say I am nervous, restless or hesitant.
\end{chapteropening}

\section[\texorpdfstring{\ja{きんちょうする} (kinchō suru)}{きんちょうする (kinchō suru)}]{\ja{きんちょうする} (kinchō suru) — to get nervous; to be tense}

\label{lesson:JA-C427-kinchou-suru}



\begin{quote}
Before the new one: say the Japanese for refreshed, then the Japanese for to endure.
\end{quote}

\subsection*{You'll want to know: \ja{きんちょうする}}
\textbf{\ja{きんちょうする}} — \emph{kinchō suru} — \textquotedblleft{}to get nervous; to be tense\textquotedblright{}.

\begin{quote}
\textbf{\ja{しあいの} \ja{まえに} \ja{きんちょうしました。}} — \emph{shiai no mae ni kinchō shimashita} — \textquotedblleft{}I got nervous before the match.\textquotedblright{}
\end{quote}

Polite \textbf{\ja{きんちょうします}}, \emph{kinchōshimasu}; a \ja{する} verb. Past (polite): \textbf{\ja{せんしゅが} \ja{きんちょうしました。}}, \emph{senshu ga kinchō shimashita}.

\begin{cousinweb}
\textbf{\ja{きんちょうする}} (\emph{kinchō suru}) is Sino-Japanese (kango) \emph{kinchō} with the native verb \emph{suru}, \textquotedblleft{}to do\textquotedblright{}: two Chinese characters read with their on'yomi, readings borrowed from Chinese. They are \emph{kin}, \textquotedblleft{}tense\textquotedblright{} (Middle Chinese \emph{kjinX}) and \emph{chō}, \textquotedblleft{}stretch\textquotedblright{} (\emph{trjang}) (KANJIDIC2; Middle Chinese in Baxter's transcription).
\end{cousinweb}

\begin{grammarlens}[title={What you've built}]
One more word for describing what you lived and felt.
\end{grammarlens}

\subsection*{Your turn}
\begin{itemize}
\raggedright
  \item \emph{Say it:} \emph{kinchō suru}
  \item \emph{Say it:} \emph{kinchō suru}, once more
  \item \emph{Say it:} \emph{shiai no mae ni kinchō shimashita} — I got nervous before the match.
  \item \emph{Recall:} say the Japanese for refreshed, then the Japanese for to endure, then say \emph{kinchō suru} again
\end{itemize}

\begin{tcolorbox}[breakable,colback=teal!4,colframe=teal!35!black,title={Before you move on}]
What does \ja{きんちょうする} mean? (\textquotedblleft{}To get nervous; to be tense\textquotedblright{}.) Say it once more.
\end{tcolorbox}
\section[\texorpdfstring{\ja{どきどきする} (dokidoki suru)}{どきどきする (dokidoki suru)}]{\ja{どきどきする} (dokidoki suru) — to have a pounding heart}

\label{lesson:JA-C427-dokidoki-suru}



\begin{quote}
Before the new one: say the Japanese for to endure, then the Japanese for to get nervous.
\end{quote}

\subsection*{You'll want to know: \ja{どきどきする}}
\textbf{\ja{どきどきする}} — \emph{dokidoki suru} — \textquotedblleft{}to have a pounding heart\textquotedblright{}.

\begin{quote}
\textbf{\ja{しけんの} \ja{まえは} \ja{むねが} \ja{どきどきします。}} — \emph{shiken no mae wa mune ga dokidoki shimasu} — \textquotedblleft{}Before an exam my heart pounds.\textquotedblright{}
\end{quote}

Polite \textbf{\ja{どきどきします}}, \emph{dokidokishimasu}; a \ja{する} verb. Past (polite): \textbf{\ja{かのじょの} \ja{むねが} \ja{どきどきしました。}}, \emph{kanojo no mune ga dokidoki shimashita}.

\begin{cousinweb}
\textbf{\ja{どきどきする}} (\emph{dokidoki suru}) is a native Japanese mimetic word, \emph{dokidoki}, inherited rather than borrowed from Chinese, made a verb with \emph{suru}, \textquotedblleft{}to do\textquotedblright{}; its older origin is not settled (JMdict).
\end{cousinweb}

\begin{grammarlens}[title={What you've built}]
One more word for describing what you lived and felt.
\end{grammarlens}

\subsection*{Your turn}
\begin{itemize}
\raggedright
  \item \emph{Say it:} \emph{dokidoki suru}
  \item \emph{Say it:} \emph{dokidoki suru}, once more
  \item \emph{Say it:} \emph{shiken no mae wa mune ga dokidoki shimasu} — Before an exam my heart pounds.
  \item \emph{Recall:} say the Japanese for to endure, then the Japanese for to get nervous, then say \emph{dokidoki suru} again
\end{itemize}

\begin{tcolorbox}[breakable,colback=teal!4,colframe=teal!35!black,title={Before you move on}]
What does \ja{どきどきする} mean? (\textquotedblleft{}To have a pounding heart\textquotedblright{}.) Say it once more.
\end{tcolorbox}
\section[\texorpdfstring{\ja{そわそわする} (sowasowa suru)}{そわそわする (sowasowa suru)}]{\ja{そわそわする} (sowasowa suru) — to be restless; to fidget}

\label{lesson:JA-C427-sowasowa-suru}



\begin{quote}
Before the new one: say the Japanese for to get nervous, then the Japanese for to have a pounding heart.
\end{quote}

\subsection*{You'll want to know: \ja{そわそわする}}
\textbf{\ja{そわそわする}} — \emph{sowasowa suru} — \textquotedblleft{}to be restless; to fidget\textquotedblright{}.

\begin{quote}
\textbf{\ja{なつやすみが} \ja{ちかいので、こどもが} \ja{そわそわしています。}} — \emph{natsuyasumi ga chikai node, kodomo ga sowasowa shite imasu} — \textquotedblleft{}The summer holidays are near, so the children are restless.\textquotedblright{}
\end{quote}

Polite \textbf{\ja{そわそわします}}, \emph{sowasowashimasu}; a \ja{する} verb. Past (polite): \textbf{\ja{こどもが} \ja{そわそわしました。}}, \emph{kodomo ga sowasowa shimashita}.

\begin{cousinweb}
\textbf{\ja{そわそわする}} (\emph{sowasowa suru}) is a native Japanese mimetic word, \emph{sowasowa}, inherited rather than borrowed from Chinese, made a verb with \emph{suru}, \textquotedblleft{}to do\textquotedblright{}; its older origin is not settled (JMdict).
\end{cousinweb}

\begin{grammarlens}[title={What you've built}]
One more word for describing what you lived and felt.
\end{grammarlens}

\subsection*{Your turn}
\begin{itemize}
\raggedright
  \item \emph{Say it:} \emph{sowasowa suru}
  \item \emph{Say it:} \emph{sowasowa suru}, once more
  \item \emph{Say it:} \emph{natsuyasumi ga chikai node, kodomo ga sowasowa shite imasu} — The summer holidays are near, so the children are restless.
  \item \emph{Recall:} say the Japanese for to get nervous, then the Japanese for to have a pounding heart, then say \emph{sowasowa suru} again
\end{itemize}

\begin{tcolorbox}[breakable,colback=teal!4,colframe=teal!35!black,title={Before you move on}]
What does \ja{そわそわする} mean? (\textquotedblleft{}To be restless; to fidget\textquotedblright{}.) Say it once more.
\end{tcolorbox}
\section[\texorpdfstring{\ja{いらいらする} (iraira suru)}{いらいらする (iraira suru)}]{\ja{いらいらする} (iraira suru) — to get irritated}

\label{lesson:JA-C427-iraira-suru}



\begin{quote}
Before the new one: say the Japanese for to have a pounding heart, then the Japanese for to be restless.
\end{quote}

\subsection*{You'll want to know: \ja{いらいらする}}
\textbf{\ja{いらいらする}} — \emph{iraira suru} — \textquotedblleft{}to get irritated\textquotedblright{}.

\begin{quote}
\textbf{\ja{でんしゃが} \ja{おくれて、いらいらしました。}} — \emph{densha ga okurete, iraira shimashita} — \textquotedblleft{}The train was late, and I got irritated.\textquotedblright{}
\end{quote}

Polite \textbf{\ja{いらいらします}}, \emph{irairashimasu}; a \ja{する} verb. Past (polite): \textbf{\ja{ちちが} \ja{いらいらしました。}}, \emph{chichi ga iraira shimashita}.

\begin{cousinweb}
\textbf{\ja{いらいらする}} (\emph{iraira suru}) is a native Japanese mimetic word, \emph{iraira}, inherited rather than borrowed from Chinese, made a verb with \emph{suru}, \textquotedblleft{}to do\textquotedblright{}; its older origin is not settled (JMdict).
\end{cousinweb}

\begin{grammarlens}[title={What you've built}]
One more word for describing what you lived and felt.
\end{grammarlens}

\subsection*{Your turn}
\begin{itemize}
\raggedright
  \item \emph{Say it:} \emph{iraira suru}
  \item \emph{Say it:} \emph{iraira suru}, once more
  \item \emph{Say it:} \emph{densha ga okurete, iraira shimashita} — The train was late, and I got irritated.
  \item \emph{Recall:} say the Japanese for to have a pounding heart, then the Japanese for to be restless, then say \emph{iraira suru} again
\end{itemize}

\begin{tcolorbox}[breakable,colback=teal!4,colframe=teal!35!black,title={Before you move on}]
What does \ja{いらいらする} mean? (\textquotedblleft{}To get irritated\textquotedblright{}.) Say it once more.
\end{tcolorbox}
\section[\texorpdfstring{\ja{ためらう} (tamerau)}{ためらう (tamerau)}]{\ja{ためらう} (tamerau) — to hesitate}

\label{lesson:JA-C427-tamerau}



\begin{quote}
Before the new one: say the Japanese for to be restless, then the Japanese for to get irritated.
\end{quote}

\subsection*{You'll want to know: \ja{ためらう}}
\textbf{\ja{ためらう}} — \emph{tamerau} — \textquotedblleft{}to hesitate\textquotedblright{}.

\begin{quote}
\textbf{\ja{かれは} \ja{へんじを} \ja{ためらいました。}} — \emph{kare wa henji o tameraimashita} — \textquotedblleft{}He hesitated to reply.\textquotedblright{}
\end{quote}

Polite \textbf{\ja{ためらいます}}, \emph{tameraimasu}; a godan verb. Past (polite): \textbf{\ja{ともだちが} \ja{すこし} \ja{ためらいました。}}, \emph{tomodachi ga sukoshi tameraimashita}.

\begin{cousinweb}
\textbf{\ja{ためらう}} (\emph{tamerau}) is a native Japanese word (wago), inherited rather than borrowed from Chinese. JMdict marks it as usually written in kana; its kanji spelling borrows two characters that each mean \textquotedblleft{}hesitate\textquotedblright{}, chosen for the meaning and not the sound (JMdict, KANJIDIC2).
\end{cousinweb}

\begin{grammarlens}[title={What you've built}]
One more word for describing what you lived and felt.
\end{grammarlens}

\subsection*{Your turn}
\begin{itemize}
\raggedright
  \item \emph{Say it:} \emph{tamerau}
  \item \emph{Say it:} \emph{tamerau}, once more
  \item \emph{Say it:} \emph{kare wa henji o tameraimashita} — He hesitated to reply.
  \item \emph{Recall:} say the Japanese for to be restless, then the Japanese for to get irritated, then say \emph{tamerau} again
\end{itemize}

\begin{tcolorbox}[breakable,colback=teal!4,colframe=teal!35!black,title={Before you move on}]
What does \ja{ためらう} mean? (\textquotedblleft{}To hesitate\textquotedblright{}.) Say it once more.
\end{tcolorbox}
